\chapter{From Prediction to Portfolio}\label{ch:ml:from-prediction-to-portfolio}

The forecast with the lowest error of the chapter's experiment produces a portfolio with a net Sharpe ratio of 0.19; a
forecast with a slightly higher error produces one of 1.52. The first was fitted on every name and learned the strong
signal in the names the desk cannot trade; the second was fitted on the names it can, where the signal is weaker and
different. A forecast is an input to a decision, and the loss it is trained on should care about what the decision does
with it. This chapter compares the three ways to make it care: train the forecast only where the portfolio uses it, train
it through the portfolio rule (\omterm{def:ml:from-prediction-to-portfolio:dfl}{decision-focused learning}), or skip the forecast and learn the portfolio rule directly
(parametric portfolio policies). It ends with the price of the second and third: they fit noise that least squares would
ignore.

\section{Why the forecast loss is not the trading loss}

\begin{definition}[Predict-then-optimise]\label{def:ml:from-prediction-to-portfolio:pto}
\emph{Predict-then-optimise}\index{predict-then-optimise} fits a forecast by a statistical loss (mean squared error, a
likelihood) and passes it to an optimiser that chooses the portfolio; the two steps are fitted separately.
\end{definition}

Book~7 (chapter~15) measured the loss between forecast and portfolio through the transfer coefficient; the loss here is
earlier, in the forecast itself. Mean squared error weighs every name and every period equally; the portfolio's profit
depends on the forecast only where the optimiser trades, in proportion to the positions it takes, and after costs. When the
forecast model is correctly specified, least squares estimates the true relation and the difference does not matter. When it
is not, the model spends its limited flexibility on the errors that matter most to the loss, which need not be those that
matter to the portfolio (Elmachtoub and Grigas, 2022).

The chapter's panel makes that concrete. Two hundred names carry four characteristics that follow slow AR(1) processes. In
the hundred names the desk cannot trade (too small, restricted), the first characteristic predicts next period's return
strongly (0.3\% per unit); in the hundred it can, only the second does, weakly (0.08\%). Returns have 3\% of noise per period.
The forecast is linear in the characteristics with one coefficient each for all names, so it cannot be right everywhere. The
portfolio rule is mean-variance on the tradeable names with a quadratic cost proxy,
\[
w_{i,t} = \frac{\hat\mu_{i,t} + 2\kappa w_{i,t-1}}{\gamma\sigma^2 + 2\kappa},
\]
with $\gamma = 20$, $\sigma = 3\%$ and $\kappa$ equal to the linear cost of 5 basis points that net returns pay on every
change in weight. Five pairs of independent panels, 1\,500 periods each, give training and test data.

Least squares on all names puts a coefficient of 1.51 (in units of 0.1\%) on the first characteristic and 0.43 on the second
(\cref{fig:ml:e2e:coefs}): it has learned mostly the untradeable signal. Out of sample its mean squared error is lower than
that of least squares on the tradeable names alone (9.007 against 9.031, in units of $10^{-4}$), and higher on the names the
desk trades (8.987 against 8.962): the better forecast by the usual measure is the worse one where it is used.

\section{Losses aligned with profit and loss}

The simplest alignment is to fit the forecast on the data the decision uses: least squares on the tradeable names finds the
second characteristic (0.83 against a true 0.80) and ignores the first. More generally, a weighted loss can weigh each name by
the size of the position the optimiser would take, or by its liquidity, and a loss on the portfolio's return rather than on the
forecast aligns the two completely.

\begin{definition}[Decision-focused learning, differentiable optimisation layer]\label{def:ml:from-prediction-to-portfolio:dfl}
\emph{Decision-focused learning}\index{decision-focused learning} trains a forecast by the quality of the decisions made with
it: the forecast is passed through the optimiser and the loss is computed on the decision's outcome (Donti, Amos and Kolter,
2017). It needs a \emph{differentiable optimisation layer}\index{differentiable optimisation layer}: an optimiser whose solution
can be differentiated with respect to its inputs, in closed form as here, by implicit differentiation of the optimality
conditions of a quadratic programme (Amos and Kolter, 2017), or for general convex problems (Agrawal and co-authors, 2019).
\end{definition}

\section{Parametric portfolio policies}

\begin{definition}[Parametric portfolio policy]\label{def:ml:from-prediction-to-portfolio:ppp}
A \emph{parametric portfolio policy}\index{parametric portfolio policy} writes the weights directly as a function of the
names' characteristics, $w_{i,t} = \theta^\top x_{i,t}/N$ around a benchmark, and chooses $\theta$ to maximise the investor's
average utility over the sample (Brandt, Santa-Clara and Valkanov, 2009); there is no forecast, and the characteristics'
covariances with returns enter only through the portfolio's outcome.
\end{definition}

The chapter trains both by gradient ascent on the Sharpe ratio of net returns over windows of 100 periods
(\cref{lst:ml:e2e:train}), starting from the all-names least-squares coefficients: the decision-focused forecast through the
mean-variance rule, the parametric \omterm{def:ml:reinforcement-learning-foundations:mdp}{policy} without its smoothing. \Cref{tab:ml:e2e:main} and \cref{fig:ml:e2e:sharpe} give the
test results.

\begin{table}[htbp]
\centering\small
\begin{tabular}{@{}p{5.6cm}cccc@{}}
\toprule
coefficients chosen by & net Sharpe & (s.d., 5 pairs) & gross Sharpe & turnover\\
\midrule
least squares, all names (predict, then optimise) & 0.19 & 0.09 & 0.60 & 2.95\\
least squares, tradeable names & 1.52 & 0.23 & 1.92 & 1.59\\
\omterm{def:ml:from-prediction-to-portfolio:dfl}{decision-focused learning} & 1.51 & 0.22 & 1.92 & 2.75\\
\omterm{def:ml:from-prediction-to-portfolio:ppp}{parametric portfolio policy} & 1.49 & 0.23 & 1.92 & 2.90\\
true coefficients & 1.57 & 0.24 & 1.97 & 1.50\\
\bottomrule
\end{tabular}
\caption{Annualised Sharpe ratios on test panels (52 periods a year) and mean turnover per period (sum of absolute weight
changes), means over five train-test pairs. Data: \texttt{ml\_e2e.results}.}\label{tab:ml:e2e:main}
\end{table}

\section{Learning through the optimiser}

Learning through the decision recovers what \omterm{def:ml:from-prediction-to-portfolio:pto}{predict-then-optimise} loses: 1.51 against 0.19, a gap of 1.33 with a standard
deviation of 0.18 over the five pairs. It learns the right direction, 1.43 on the second characteristic and almost nothing on the
first; the scale is larger than the truth's because a Sharpe ratio does not depend on it, and the turnover is higher for the
same reason. But it does no better than least squares on the tradeable names, which is simpler, faster and turns over less.
The gain came from aligning the data with the decision, not from differentiating through the optimiser; here the optimiser
is simple enough that the two coincide. \omterm{def:ml:from-prediction-to-portfolio:dfl}{Decision-focused learning} earns its complexity when the decision's sensitivity to the
forecast varies in ways a weight cannot express: constraints that bind for some names and not others, costs that depend on the
trade, risk that couples the names.

\begin{omfigure}
\begin{tikzpicture}
\begin{axis}[width=12cm,height=6cm,ybar,bar width=10pt,ymin=0,ymax=2.1,ylabel={net Sharpe ratio (test)},area legend,
  xtick={0,1,2,3,4},xticklabels={{predict,\\optimise},{least squares,\\tradeable},{decision-\\focused},{parametric\\policy},{true}},
  xticklabel style={align=center},
  tick label style={font=\small},label style={font=\small},xmin=-0.5,xmax=4.5,grid=major,grid style={gray!25},
  legend columns=2,legend style={font=\footnotesize,at={(0.5,-0.3)},anchor=north,draw=none,fill=none},table/col sep=comma]
\addplot[fill=omDef,draw=none,bar shift=-5.5pt,error bars/.cd,y dir=both,y explicit]
  table[x=i,y=net,y error=sd]{figdata/ml/22-from-prediction-to-portfolio/sharpe.csv};
\addplot[fill=omThm,draw=none,bar shift=5.5pt,error bars/.cd,y dir=both,y explicit]
  table[x=i,y=netover,y error=sdover]{figdata/ml/22-from-prediction-to-portfolio/sharpe.csv};
\legend{{1\,500 periods, 4 characteristics},{300 periods, 20 characteristics}}
\end{axis}
\end{tikzpicture}

\omcaption{Net Sharpe ratio on test panels by the way the forecast's coefficients were chosen, mean and standard deviation over
five train-test pairs, with ample data and with little data and many useless characteristics. Data:
\texttt{ml\_e2e.summary}.}\label{fig:ml:e2e:sharpe}
\end{omfigure}

\begin{omfigure}
\begin{tikzpicture}
\begin{axis}[width=11cm,height=5.6cm,ybar,bar width=10pt,ylabel={coefficient (units of 0.1\%)},area legend,
  xtick={0,1,2,3,4},xticklabels={{predict,\\optimise},{least squares,\\tradeable},{decision-\\focused},{parametric\\policy},{true}},
  xticklabel style={align=center},
  tick label style={font=\small},label style={font=\small},xmin=-0.5,xmax=4.5,grid=major,grid style={gray!25},
  legend columns=2,legend style={font=\footnotesize,at={(0.5,-0.3)},anchor=north,draw=none,fill=none},table/col sep=comma]
\addplot[fill=gray,draw=none,bar shift=-5.5pt] table[x=i,y=b1]{figdata/ml/22-from-prediction-to-portfolio/coefs.csv};
\addplot[fill=omDef,draw=none,bar shift=5.5pt] table[x=i,y=b2]{figdata/ml/22-from-prediction-to-portfolio/coefs.csv};
\legend{first characteristic (untradeable signal),second characteristic (tradeable signal)}
\end{axis}
\end{tikzpicture}

\omcaption{The forecast's coefficients on the two informative characteristics, means over five training panels. Data:
\texttt{ml\_e2e.results}.}\label{fig:ml:e2e:coefs}
\end{omfigure}

\section{When end-to-end learning overfits}

A loss computed on the portfolio's net Sharpe ratio is a noisy, low-dimensional summary of the data; fitting many parameters to
it is fitting to noise. With 300 training periods and sixteen useless characteristics added, least squares on the tradeable
names still reaches 1.04 out of sample, \omterm{def:ml:from-prediction-to-portfolio:dfl}{decision-focused learning} 0.89 and the parametric \omterm{def:ml:reinforcement-learning-foundations:mdp}{policy} 0.87 (\cref{fig:ml:e2e:sharpe}),
with higher turnover (3.49 against 1.88): the end-to-end fits spent their freedom on characteristics that happened to pay in
the training sample. Even with 1\,500 periods, the decision-focused coefficients on the two useless characteristics of the
first panel were $-0.44$ and 0.20, against $-0.22$ and 0.13 for least squares on the tradeable names. The remedies are those of any small-sample
fit: fewer parameters, penalties, \omterm{def:ml:trees-and-boosting:early}{early stopping} on a validation block, and a least-squares starting point that the
end-to-end fit must beat on held-out data.

\begin{method}[From forecast to portfolio]\label{met:ml:e2e:method}
\begin{enumerate}
\item Fit forecasts on the names, periods and horizons the portfolio uses, weighted by how much the portfolio depends on them.
\item Evaluate forecasts by the net performance of the portfolio built from them, never by forecast error alone.
\item Use \omterm{def:ml:from-prediction-to-portfolio:dfl}{decision-focused learning} or parametric policies when the decision's sensitivity to the forecast cannot be expressed
as a weight, with strong regularisation and a least-squares benchmark.
\item Report the spread across independent samples; a gap smaller than it is no gap.
\end{enumerate}
\end{method}

\section{Tutorial: the loss that pays}\label{tut:ml:from-prediction-to-portfolio:tut}

\begin{tutorial}
\textbf{Goal.} Compare \omterm{def:ml:from-prediction-to-portfolio:pto}{predict-then-optimise}, aligned least squares, \omterm{def:ml:from-prediction-to-portfolio:dfl}{decision-focused learning} and a parametric \omterm{def:ml:reinforcement-learning-foundations:mdp}{policy} on five
train-test pairs, then with little data and many characteristics. \textbf{End state:} \cref{tab:ml:e2e:main},
\cref{fig:ml:e2e:sharpe,fig:ml:e2e:coefs}.
\begin{tutsteps}
\item \textbf{The portfolio rule and its evaluation.}
\omcode{code/firm/e2eport/firm_e2eport.py}{45}{68}{A differentiable mean-variance rule with a cost proxy, and the net-return evaluation.}\label{lst:ml:e2e:run}
\item \textbf{Learning through the decision.}
\omcode{code/firm/e2eport/firm_e2eport.py}{76}{100}{Gradient ascent on the net Sharpe ratio through the rule.}\label{lst:ml:e2e:train}
\item \textbf{Run} \texttt{ml\_e2e.results()}, \texttt{overfit()}, \texttt{summary()}, \texttt{forecast\_errors()} and
\texttt{fig\_e2e.py} (about 45 seconds on one core).
\end{tutsteps}
\textbf{What to change next.} Add a name-level position limit that binds for some names and see whether aligned least squares
can still match \omterm{def:ml:from-prediction-to-portfolio:dfl}{decision-focused learning}; add \omterm{def:ml:trees-and-boosting:early}{early stopping} on a validation block to the end-to-end fit.
\end{tutorial}

\section{Build: end-to-end portfolios}\label{bld:ml:from-prediction-to-portfolio:e2eport}

\begin{build}
\bfield{Purpose} Forecasts judged and, where it pays, trained by the portfolios they produce.
\bfield{Interface} \texttt{char\_panel(seed, T, N, K, n\_noise, phi)}, \texttt{mv\_layer(mu, w\_prev, kappa, gamma, sigma)},
\texttt{run(b, X, R, tradeable, cost, gamma, sigma, policy)}, \texttt{least\_squares(X, R, rows)},
\texttt{train\_decision(X, R, tradeable, init, policy, epochs, lr, window, seed)}.
\bfield{Rules} Every end-to-end result is reported against aligned least squares on the same held-out data, with its spread over
independent samples.
\bfield{Acceptance tests} \texttt{code/firm/e2eport/tests/}: the closed-form layer equals Book~7's \texttt{firm.portcons}
solution when costs are zero; the layer's gradient matches finite differences; the evaluation charges costs on weight changes;
decision-focused training on a clean problem recovers the direction of the true coefficients.
\bfield{Stretch} A general quadratic-programme layer by implicit differentiation, with a factor risk model; name-level limits;
\omterm{def:ml:trees-and-boosting:early}{early stopping}.
\end{build}

\begin{omsources}
\item M.~W.~Brandt, P.~Santa-Clara and R.~Valkanov, ``Parametric portfolio policies: exploiting characteristics in the
cross-section of equity returns'', \emph{Review of Financial Studies} 22(9), 2009.
\item A.~N.~Elmachtoub and P.~Grigas, ``Smart `predict, then optimize''', \emph{Management Science} 68(1), 2022.
\item P.~L.~Donti, B.~Amos and J.~Z.~Kolter, ``Task-based end-to-end model learning in stochastic optimization'',
arXiv:1703.04529, 2017.
\item B.~Amos and J.~Z.~Kolter, ``OptNet: differentiable optimization as a layer in neural networks'', arXiv:1703.00443, 2017.
\item A.~Agrawal and co-authors, ``Differentiable convex optimization layers'', arXiv:1910.12430, 2019.
\item A.~Butler and R.~H.~Kwon, ``Integrating prediction in mean-variance portfolio optimization'', \emph{Quantitative Finance}
23(3), 2023.
\end{omsources}

\section{Exercises}

\begin{exercise}[$\star$]\label{exo:ml:from-prediction-to-portfolio:1}
Derive the rule $w = (\mu + 2\kappa w_{\mathrm{prev}})/(\gamma\sigma^2 + 2\kappa)$ from the one-name objective. What does it
become with $\kappa = 0$ and as $\kappa\to\infty$?
\end{exercise}

\begin{exercise}[$\star$]\label{exo:ml:from-prediction-to-portfolio:2}
With the chapter's parameters, what fraction of the gap to its target does a name's weight close each period?
\end{exercise}

\begin{exercise}[$\star$]\label{exo:ml:from-prediction-to-portfolio:3}
Why does a Sharpe-ratio loss leave the scale of the coefficients undetermined, and what does that do to turnover?
\end{exercise}

\begin{exercise}[$\star\star$]\label{exo:ml:from-prediction-to-portfolio:4}
Explain why the all-names least-squares forecast has the lower error overall and the higher error on the tradeable names.
\end{exercise}

\begin{exercise}[$\star\star$]\label{exo:ml:from-prediction-to-portfolio:5}
When would \omterm{def:ml:from-prediction-to-portfolio:dfl}{decision-focused learning} beat least squares on the tradeable names in this setting?
\end{exercise}

\begin{exercise}[$\star\star$]\label{exo:ml:from-prediction-to-portfolio:6}
\emph{Find the flaw.} ``We trained the network end to end on the portfolio's Sharpe ratio and got 3.1 in the backtest, against
1.2 for the forecast trained on returns, so end-to-end training is worth it.''
\end{exercise}

\begin{exercise}[$\star\star\star$]\label{exo:ml:from-prediction-to-portfolio:7}
\emph{Coding.} Add a ridge penalty to the decision-focused training in the small-sample scenario and choose its strength on a
validation block. Does it close the gap to aligned least squares?
\end{exercise}

\begin{exercise}[$\star\star\star$]\label{exo:ml:from-prediction-to-portfolio:8}
For the unconstrained mean-variance problem $\max_w\mu^\top w - \frac\gamma2w^\top\Sigma w$, compute $\partial w^*/\partial\mu$
and explain how the gradient of a portfolio loss flows back to the forecast.
\end{exercise}

\section{Problem: The Loss That Pays}

\begin{problem}[{Weekend problem --- train on what you trade}]\label{pb:ml:from-prediction-to-portfolio:1}
The chapter's panel, portfolio rule and five ways to choose the forecast.

\medskip\noindent\textbf{Part I --- The setup.}
\begin{enumerate}
\item What does each half of the universe contain, and what can the linear forecast represent?
\item What does the portfolio rule do, and how are costs charged?
\item What coefficients does least squares on all names learn, and why?
\item How do the two least-squares forecasts compare by mean squared error?
\end{enumerate}

\medskip\noindent\textbf{Part II --- The comparison.}
\begin{enumerate}[resume]
\item Give the net and gross Sharpe ratios and the turnover of the five methods.
\item What is the gap between \omterm{def:ml:from-prediction-to-portfolio:dfl}{decision-focused learning} and \omterm{def:ml:from-prediction-to-portfolio:pto}{predict-then-optimise}, and its spread?
\item Why do aligned least squares and \omterm{def:ml:from-prediction-to-portfolio:dfl}{decision-focused learning} tie?
\item Why is the parametric \omterm{def:ml:reinforcement-learning-foundations:mdp}{policy}'s turnover highest?
\end{enumerate}

\medskip\noindent\textbf{Part III --- \omterm{def:ml:why-financial-machine-learning-is-different:generalisation}{Overfitting}.}
\begin{enumerate}[resume]
\item What happens with 300 periods and twenty characteristics?
\item Which coefficients show the overfit?
\item What regularisation would you use?
\item How would you pick between the methods on real data?
\end{enumerate}

\medskip\noindent\textbf{Part IV --- The verdict.}
\begin{enumerate}[resume]
\item State the \emph{named result}: the net Sharpe ratio of \omterm{def:ml:from-prediction-to-portfolio:pto}{predict-then-optimise} against \omterm{def:ml:from-prediction-to-portfolio:dfl}{decision-focused learning} and the
parametric \omterm{def:ml:reinforcement-learning-foundations:mdp}{policy}, with the gap's spread over seeds.
\item When is \omterm{def:ml:from-prediction-to-portfolio:pto}{predict-then-optimise} enough?
\item What would make \omterm{def:ml:from-prediction-to-portfolio:dfl}{decision-focused learning} worth its cost?
\item How does this relate to the transfer coefficient?
\item What should a research report on a new forecast contain, after this chapter?
\item Where do costs enter the training loss, and where should they?
\item How would you do this with a neural forecast?
\item In one sentence: what should a forecast be trained to do?
\end{enumerate}
\end{problem}

\section{Interview questions}

\begin{interviewq}[$\star$ \iqroles{researcher}]\label{iq:ml:from-prediction-to-portfolio:1}
Why can a more accurate forecast produce a worse portfolio?
\end{interviewq}

\begin{interviewq}[$\star\star$ \iqroles{researcher, mle}]\label{iq:ml:from-prediction-to-portfolio:2}
What is \omterm{def:ml:from-prediction-to-portfolio:dfl}{decision-focused learning}, and how do you differentiate through an optimiser?
\end{interviewq}

\begin{interviewq}[$\star\star$ \iqroles{researcher}]\label{iq:ml:from-prediction-to-portfolio:3}
Describe Brandt, Santa-Clara and Valkanov's parametric portfolio policies. What are their advantages and risks?
\end{interviewq}

\begin{interviewq}[$\star\star$ \iqroles{researcher, trader}]\label{iq:ml:from-prediction-to-portfolio:4}
How would you include transaction costs in the training of an alpha model?
\end{interviewq}

\begin{interviewq}[$\star\star$ \iqroles{mle}]\label{iq:ml:from-prediction-to-portfolio:5}
An end-to-end portfolio network backtests far better than \omterm{def:ml:from-prediction-to-portfolio:pto}{predict-then-optimise}. What do you check?
\end{interviewq}

\begin{interviewq}[$\star\star\star$ \iqroles{researcher}]\label{iq:ml:from-prediction-to-portfolio:6}
Show how the KKT conditions of a quadratic programme give the derivative of its solution with respect to the linear term.
\end{interviewq}
